\documentclass{article}

% This file follows the ICML 2026 preprint template used by the arXiv source of:
% "FineInstructions: Scaling Synthetic Instructions to Pre-Training Scale"
% https://arxiv.org/abs/2601.22146
%
% Overleaf setup:
% 1. Create a project from the ICML 2026 template, or upload the arXiv TeX
%    source files for 2601.22146.
% 2. Ensure icml2026.sty and icml2026.bst are in the project root.
% 3. Set this file as the main TeX file.

\usepackage[utf8]{inputenc}
\usepackage[T1]{fontenc}
\usepackage[brazilian]{babel}

\usepackage{microtype}
\usepackage{graphicx}
\usepackage{subcaption}
\usepackage{booktabs}
\usepackage{multirow}
\usepackage{enumitem}
\usepackage{xcolor}
\usepackage{tikz}
\usepackage{url}
\usepackage{hyperref}

\usepackage[preprint]{icml2026}

\usepackage{amsmath}
\usepackage{amssymb}
\usepackage{mathtools}
\usepackage{amsthm}
\usepackage[capitalize,noabbrev]{cleveref}

\usetikzlibrary{arrows.meta,positioning,shapes.geometric}

\setlist[itemize]{leftmargin=*,topsep=0.2em,itemsep=0.1em}
\setlist[enumerate]{leftmargin=*,topsep=0.2em,itemsep=0.1em}

\newcommand{\datasetname}{\textsc{Wiki-Brazil-SFT}}
\newcommand{\toolname}{\texttt{sft-dataset-creator}}
\newcommand{\runname}{\texttt{wiki-brazil-all-tasks-22596}}

\icmltitlerunning{Pipeline Deterministica para Dados Sinteticos de SFT}

\begin{document}

\twocolumn[
  \icmltitle{Uma Pipeline Deterministica para Geracao e Filtragem de Dados Sinteticos de SFT em Portugues Brasileiro}

  \begin{icmlauthorlist}
    \icmlauthor{Autor(es) a definir}{inst}
  \end{icmlauthorlist}

  \icmlaffiliation{inst}{Instituicao a definir}
  \icmlcorrespondingauthor{Autor correspondente}{email@dominio.org}
  \icmlkeywords{Synthetic Data, Instruction Tuning, SFT, Deterministic Evaluation, Portuguese}

  \vskip 0.3in
]

\printAffiliationsAndNotice{}

\begin{abstract}
Dados supervisionados de instrucao e resposta sao um recurso central para
ajustar modelos de linguagem ao uso interativo, mas sua construcao manual e
cara e dificil de escalar. Este artigo descreve uma pipeline reprodutivel para
gerar dados sinteticos de supervised fine-tuning (SFT) em portugues brasileiro
a partir de artigos da Wikipedia. O metodo combina selecao e chunking
deterministicos do corpus, planejamento balanceado de tipos de tarefa e
dificuldade, geracao batelada com um unico modelo instrucional local e uma
etapa de filtragem inteiramente deterministica. Diferentemente de pipelines que
dependem de um segundo modelo julgador, a avaliacao principal aceita apenas
candidatos que satisfazem regras verificaveis de formato, auto-contencao,
nao duplicacao e ancoragem em spans de evidencia extraidos do documento de
origem. Em uma execucao sobre o corpus \texttt{costadev00/wiki-brazil}, o
sistema planejou 22.596 exemplos a partir de 1.614 documentos elegiveis,
gerou 35.876 tentativas e exportou 20.378 exemplos aceitos em tres formatos
de SFT. A analise do ledger SQLite da execucao mostra taxa de aceitacao de
56,8\%, com rejeicoes dominadas por problemas de evidencia, duplicatas
normalizadas e referencias ao contexto oculto.
\end{abstract}

\noindent
\textbf{Palavras-chave:} dados sinteticos; SFT; instruction tuning; avaliacao
deterministica; grounding; portugues brasileiro.

\section{Introducao}

Modelos de linguagem de grande escala costumam ser treinados primeiro em
objetivos auto-supervisionados sobre grandes colecoes de texto e, depois,
adaptados com exemplos supervisionados de instrucoes e respostas
\citep{ouyang2022training,wei2022finetuned}. Esse segundo estagio, usualmente
chamado de instruction tuning ou supervised fine-tuning (SFT), aproxima o
comportamento do modelo do modo como usuarios formulam pedidos em sistemas
interativos. Entretanto, datasets supervisionados de alta qualidade sao caros,
limitados em escala e frequentemente concentrados em ingles.

Uma alternativa e transformar documentos existentes em pares sinteticos de
instrucao e resposta. Trabalhos recentes investigam geracao de dados
instrucionais a partir de queries, templates, documentos e modelos auxiliares,
incluindo pipelines em escala de pre-treinamento como FineInstructions
\citep{patel2026fineinstructions}. O presente trabalho parte de motivacao
semelhante, mas foca em um cenario menor e auditavel: a construcao de um
dataset sintetico em portugues brasileiro, com rastreabilidade por documento e
filtragem deterministica das geracoes.

Descrevemos a implementacao de \toolname{}, uma biblioteca e CLI para planejar,
gerar, avaliar e exportar datasets sinteticos de SFT. A contribuicao central
deste artigo e a formalizacao da pipeline de filtragem: um conjunto de gates
automaticos que separa a tarefa generativa, realizada por um LLM, da decisao de
aceitacao, realizada por regras verificaveis. Essa separacao reduz dependencia
de julgamento opaco por modelos adicionais e torna cada decisao inspecionavel a
partir dos artefatos da execucao.

\paragraph{Contribuicoes.}
Este artigo apresenta:
(i) uma pipeline reprodutivel para transformar artigos da Wikipedia em exemplos
de SFT em portugues;
(ii) um planejamento deterministico de slots por documento, tarefa e
dificuldade;
(iii) uma estrategia de filtragem baseada em schema, auto-contencao,
deduplicacao e evidencia;
e (iv) uma analise quantitativa de uma run com 35.876 tentativas e 20.378
exemplos aceitos.

\begin{figure*}[t]
\centering
\begin{tikzpicture}[
  node distance=0.85cm and 0.9cm,
  box/.style={draw, rounded corners=2pt, align=center, minimum height=0.75cm, text width=2.65cm, font=\small},
  gate/.style={draw, diamond, aspect=1.8, align=center, text width=2.2cm, font=\small},
  arrow/.style={-Latex, thick}
]
\node[box] (source) {Fonte\\Hugging Face};
\node[box, right=of source] (profile) {Filtro\\\texttt{wikipedia\_ptbr}};
\node[box, right=of profile] (chunk) {Chunking\\8.000 chars};
\node[box, right=of chunk] (plan) {Plano\\deterministico};
\node[box, below=of plan] (gen) {Geracao\\vLLM local};
\node[gate, left=of gen] (schema) {JSON e schema\\validos?};
\node[gate, left=of schema] (quality) {Gates\\deterministicos};
\node[box, left=of quality] (accepted) {Exemplo\\aceito};
\node[box, below=of quality] (retry) {Nova tentativa\\ate limite};
\node[box, below=of accepted] (export) {Exportacao\\com trava final};

\draw[arrow] (source) -- (profile);
\draw[arrow] (profile) -- (chunk);
\draw[arrow] (chunk) -- (plan);
\draw[arrow] (plan) -- (gen);
\draw[arrow] (gen) -- (schema);
\draw[arrow] (schema) -- node[above,font=\scriptsize] {sim} (quality);
\draw[arrow] (quality) -- node[above,font=\scriptsize] {accept} (accepted);
\draw[arrow] (accepted) -- (export);
\draw[arrow] (schema.south) |- node[pos=0.25,right,font=\scriptsize] {nao} (retry.east);
\draw[arrow] (quality.south) -- node[right,font=\scriptsize] {reject} (retry.north);
\draw[arrow] (retry.east) -| (gen.south);
\end{tikzpicture}
\caption{Visao geral da pipeline. O LLM propoe candidatos estruturados, mas a
aceitacao depende de validacoes deterministicas e de uma trava adicional antes
da exportacao.}
\label{fig:pipeline}
\end{figure*}

\section{Corpus e Pre-processamento}

A fonte usada na execucao principal foi o dataset Hugging Face
\texttt{costadev00/wiki-brazil}, split \texttt{train}, em modo streaming. O
campo de identificacao foi \texttt{\_\_row\_index\_\_}; texto, titulo, secoes e
licenca foram mapeados dos campos \texttt{text}, \texttt{title},
\texttt{sections} e \texttt{license}.

Antes de qualquer chamada ao modelo, cada registro foi normalizado como um
documento interno e passou pelo perfil \texttt{wikipedia\_ptbr}. Esse perfil
remove documentos vazios, paginas fora do namespace principal, categorias,
templates, arquivos, paginas de ajuda, portais, redirects, paginas de
desambiguacao e textos com menos de 300 caracteres. Essa filtragem inicial
reduz o risco de gerar exemplos a partir de paginas administrativas, muito
curtas ou pouco informativas.

Os documentos selecionados foram divididos em chunks de ate 8.000 caracteres,
com overlap de 400 caracteres. Cada chunk funciona como uma secao de evidencia
e como contexto de entrada para a geracao. O chunking tem tres funcoes:
controlar o tamanho de contexto, permitir offsets de evidencia verificaveis e
distribuir tentativas de um mesmo documento por regioes diferentes do artigo.

\section{Planejamento Deterministico}

O planejamento nao chama o modelo. Ele apenas fixa, de forma reprodutivel, a
relacao entre documentos, tarefas, dificuldades e chunks. A run analisada,
\runname{}, utilizou seed 42 e selecionou todos os 1.614 documentos elegiveis
(\texttt{fraction=1.0}). O alvo foi de 22.596 exemplos, com minimo e maximo de
14 exemplos por documento. Como $22\,596 = 1\,614 \times 14$, cada documento
recebeu 14 slots planejados.

As tarefas foram balanceadas globalmente entre 14 tipos:
\texttt{classification}, \texttt{closed\_qa}, \texttt{comparison},
\texttt{concept\_explanation}, \texttt{definition},
\texttt{didactic\_explanation}, \texttt{fact\_checking},
\texttt{information\_extraction}, \texttt{rewrite}, \texttt{short\_answer},
\texttt{structured\_extraction}, \texttt{summarization}, \texttt{taxonomy} e
\texttt{timeline}. Cada tipo recebeu exatamente 1.614 slots. As dificuldades
seguiram pesos 0,25 para \texttt{easy}, 0,50 para \texttt{medium} e 0,25 para
\texttt{hard}. A atribuicao documento--tarefa foi embaralhada pela seed; logo,
o balanceamento e global, nao uma garantia de uma tarefa de cada tipo por
documento.

\begin{table}[t]
\centering
\caption{Configuracao principal da run \runname{}.}
\label{tab:config}
\small
\begin{tabular}{ll}
\toprule
Parametro & Valor \\
\midrule
Documentos elegiveis & 1.614 \\
Slots planejados & 22.596 \\
Exemplos por documento & 14 min.; 14 max. \\
Tarefas & 14 tipos balanceados \\
Dificuldades & 25\% easy; 50\% medium; 25\% hard \\
Tentativas por slot & ate 5 \\
Reserva de documentos & 0\% \\
Chunk & 8.000 chars; overlap 400 \\
Seed & 42 \\
\bottomrule
\end{tabular}
\end{table}

\section{Geracao dos Candidatos}

A geracao foi realizada com um unico modelo local, executado por
\texttt{vllm\_local}: \texttt{google/gemma-4-31B-it-qat-w4a16-ct}. A
configuracao usou \texttt{tensor\_parallel\_size=4},
\texttt{quantization=compressed-tensors}, \texttt{kv\_cache\_dtype=fp8},
\texttt{max\_num\_batched\_tokens=16384}, temperatura 0,1, janela de contexto
de 65.536 tokens, limite de entrada de 53.248 tokens, limite de saida de 4.096
tokens e ate 32 requisicoes em voo. O uso de vLLM permite manter o gerador
carregado entre rodadas e executar requisicoes em lote \citep{kwon2023vllm}.

Cada requisicao contem um prompt de sistema fixo e um payload JSON com tipo de
tarefa, lingua, dificuldade, identificador do documento, titulo, indicador de
truncamento e trecho-fonte. O modelo deve retornar JSON compativel com um
schema fechado contendo \texttt{instruction}, \texttt{input}, \texttt{output},
\texttt{risk} e ao menos um item de \texttt{evidence}, com \texttt{section\_id},
\texttt{start}, \texttt{end} e, opcionalmente, \texttt{quote}.

O prompt instrui o modelo a usar apenas fatos do trecho e a nao mencionar texto,
documento, contexto, passagem ou fonte oculta. Essa instrucao, contudo, e
tratada apenas como uma orientacao de geracao: a conformidade e verificada por
filtros deterministas na etapa seguinte.

\section{Filtragem Deterministica}

A filtragem acontece em camadas. Um candidato so chega ao dataset final se
passar por todas as validacoes descritas nesta secao.

\subsection{Parse e validacao estrutural}

Primeiro, a resposta precisa ser parseavel como JSON e respeitar o schema
esperado. Se o backend retorna JSON invalido, campos ausentes ou spans
impossiveis de construir como \texttt{EvidenceSpan}, a tentativa e registrada
como \texttt{error}, nao como rejeicao semantica. Na run principal, houve 3.180
erros desse tipo: 3.157 associados a parse ou schema JSON e 23 a validacao
Pydantic.

\subsection{Campos essenciais}

Para candidatos estruturalmente validos, o avaliador verifica se
\texttt{instruction} e \texttt{output} nao estao vazios e se a instrucao tem
ao menos 12 caracteres. Falhas recebem as issues
\texttt{empty\_instruction\_or\_output} ou
\texttt{instruction\_too\_short}.

\subsection{Auto-contencao}

O exemplo final nao deve depender de um contexto invisivel para o usuario. Por
isso, a pipeline rejeita candidatos que contenham expressoes como ``de acordo
com o texto'', ``com base no documento'', ``according to the passage'' ou
equivalentes em portugues, ingles e espanhol. A deteccao e aplicada aos campos
visiveis ao usuario: \texttt{instruction}, \texttt{input} e \texttt{output}.
A issue associada e \texttt{source\_reference\_in\_candidate}.

\subsection{Deduplicacao}

Cada candidato recebe uma impressao digital textual obtida por normalizacao
Unicode, remocao de marcas combinantes, \texttt{casefold}, colapso de espacos
e concatenacao de \texttt{instruction}, \texttt{input} e \texttt{output}. Se a
impressao digital ja existir entre exemplos aceitos, o candidato e rejeitado
com \texttt{duplicate\_candidate}.

\subsection{Evidencia}

Cada candidato precisa apontar para ao menos um span de evidencia em uma secao
do documento usado na tentativa. Para cada span, a pipeline verifica se
\texttt{section\_id} existe, se \texttt{start} e \texttt{end} estao dentro dos
limites da secao, se \texttt{end > start}, se o \texttt{quote} fornecido bate
com o trecho correspondente e se o span recuperado tem ao menos 20 caracteres.

As issues possiveis sao \texttt{invalid\_section},
\texttt{invalid\_evidence\_offsets}, \texttt{evidence\_quote\_mismatch},
\texttt{missing\_grounding} e \texttt{weak\_grounding}. A ultima e registrada
como diagnostico, mas nao causa rejeicao isoladamente. As demais sao criticas.

\subsection{Decisao e novas tentativas}

O avaliador acumula as issues do candidato. Se qualquer issue critica aparece,
o veredito e \texttt{reject}; caso contrario, e \texttt{accept}. Issues
criticas incluem campos vazios, instrucao curta, secao invalida, offsets
invalidos, quote divergente, ausencia de grounding, referencia a fonte oculta e
duplicata normalizada.

Slots nao aceitos permanecem pendentes ate o limite de tentativas. Na run
analisada, cada slot podia ser tentado ate cinco vezes. Como
\texttt{reserve\_fraction=0.0}, as novas tentativas continuaram usando os
documentos planejados originalmente.

\subsection{Trava final de exportacao}

Antes da escrita dos arquivos finais, o exportador reabre os candidatos aceitos
no banco SQLite e executa novamente a deteccao de referencia a fonte oculta.
Se qualquer exemplo aceito ainda depender do texto-fonte invisivel, a
exportacao falha. Essa etapa funciona como uma garantia de ultima linha.

\section{Resultados}

A Tabela~\ref{tab:results} resume a execucao principal. A run terminou como
\texttt{partial}, pois 2.218 slots chegaram ao limite de cinco tentativas sem
produzir candidato aceito. Ainda assim, os arquivos exportados contem 20.378
exemplos em cada uma das tres visoes: \texttt{messages},
\texttt{prompt\_completion} e \texttt{alpaca}. Como os splits configurados
foram \texttt{train=1.0}, \texttt{validation=0.0} e \texttt{test=0.0}, apenas
o split de treino foi produzido.

\begin{table}[t]
\centering
\caption{Resumo quantitativo da run principal.}
\label{tab:results}
\small
\begin{tabular}{lr}
\toprule
Metrica & Valor \\
\midrule
Alvo planejado & 22.596 \\
Exemplos aceitos & 20.378 \\
Tentativas totais & 35.876 \\
Rejeicoes avaliadas & 12.318 \\
Erros de parse/schema/backend & 3.180 \\
Slots pendentes & 2.218 \\
Taxa de aceitacao & 56,8\% \\
Exemplos julgados por LLM & 0 \\
\bottomrule
\end{tabular}
\end{table}

\begin{table*}[t]
\centering
\caption{Principais issues registradas nos candidatos rejeitados. As contagens
nao sao mutuamente exclusivas: um candidato pode acumular mais de uma issue.}
\label{tab:issues}
\small
\begin{tabular}{lrl}
\toprule
Issue & Ocorrencias & Interpretacao \\
\midrule
\texttt{missing\_grounding} & 7.925 & Nenhum span recuperavel sustentou o candidato. \\
\texttt{invalid\_section} & 7.890 & O \texttt{section\_id} indicado nao existia no chunk usado. \\
\texttt{duplicate\_candidate} & 2.870 & Conteudo repetido apos normalizacao. \\
\texttt{source\_reference\_in\_candidate} & 1.855 & Instrucao, input ou output dependiam do texto oculto. \\
\texttt{weak\_grounding} & 927 & Evidencia recuperada era curta demais. \\
\texttt{invalid\_evidence\_offsets} & 232 & Offsets fora dos limites ou inconsistentes. \\
\texttt{empty\_instruction\_or\_output} & 16 & Instrucao ou resposta vazia. \\
\texttt{evidence\_quote\_mismatch} & 9 & Quote fornecido nao batia com o span. \\
\bottomrule
\end{tabular}
\end{table*}

A Tabela~\ref{tab:attempts} mostra que a maior parte dos exemplos aceitos veio
da primeira tentativa. Ainda assim, as tentativas adicionais recuperaram 2.610
exemplos que seriam perdidos em uma pipeline com geracao unica por slot.

\begin{table}[t]
\centering
\caption{Exemplos aceitos por numero da tentativa.}
\label{tab:attempts}
\small
\begin{tabular}{lr}
\toprule
Tentativa & Aceitos \\
\midrule
1 & 17.768 \\
2 & 1.499 \\
3 & 616 \\
4 & 303 \\
5 & 192 \\
\bottomrule
\end{tabular}
\end{table}

O par mais frequente de issues foi \texttt{invalid\_section} +
\texttt{missing\_grounding}, com 7.180 ocorrencias. Esse resultado indica que
uma parcela relevante das rejeicoes decorreu de referencias de evidencia
mal alinhadas ao chunk efetivamente usado, e nao necessariamente de baixa
qualidade textual da resposta. Em outras palavras, a principal fragilidade
observada esta na capacidade do gerador de apontar corretamente para a
evidencia, nao apenas em produzir respostas plausiveis.

\section{Discussao}

Os resultados mostram que a separacao entre geracao e aceitacao oferece uma
forma simples de auditar datasets sinteticos. O modelo e usado para propor
exemplos, mas nao para decidir se eles entram no corpus final. Essa decisao e
tomada por regras transparentes: parse, schema, auto-contencao, deduplicacao e
evidencia. A consequencia pratica e que cada exemplo aceito pode ser rastreado
ate seu documento, chunk, tarefa, tentativa e avaliacao.

O desenho tambem evita uma dependencia operacional importante: nao ha segundo
modelo de avaliacao na run analisada. O relatorio final registra
\texttt{llm\_judged\_examples=0}, \texttt{llm\_judge\_coverage=0.0} e
\texttt{semantic\_judge\_configured=false}. Isso reduz custo e variancia, mas
tambem limita o tipo de qualidade que pode ser garantida automaticamente.

\section{Limitacoes}

A avaliacao deterministica verifica forma, rastreabilidade e sinais de
grounding, mas nao substitui uma checagem semantica completa de todas as
afirmacoes da resposta. Um span valido demonstra que existe evidencia
recuperavel, mas nao garante que toda a resposta seja integralmente sustentada
pelo trecho.

Outra limitacao e que o filtro de duplicatas e exato apos normalizacao textual.
Ele detecta repeticoes triviais, mas nao necessariamente parafrases
semanticamente equivalentes. Alem disso, como a run principal terminou como
\texttt{partial}, o corpus exportado contem 20.378 exemplos aceitos, nao os
22.596 planejados.

Por fim, o volume de \texttt{invalid\_section} sugere que a relacao entre
prompt, chunk selecionado e IDs de secao pode ser melhorada. Trabalhos futuros
podem explorar representacoes de evidencia mais faceis para o modelo, validacao
incremental durante a decodificacao ou geracao em duas etapas: primeiro o
exemplo, depois a atribuicao de evidencia.

\section{Conclusao}

Este artigo formalizou uma pipeline para gerar e filtrar dados sinteticos de
SFT em portugues brasileiro a partir de documentos da Wikipedia. A execucao
analisada mostra que e possivel produzir dezenas de milhares de tentativas com
um unico modelo local e exportar mais de 20 mil exemplos aceitos usando apenas
gates deterministas. A principal contribuicao metodologica e tornar explicita a
fronteira entre geracao e validacao: o LLM produz candidatos, mas a entrada no
dataset depende de propriedades verificaveis e auditaveis.

\section*{Disponibilidade dos artefatos}

Os artefatos locais da execucao incluem \texttt{config.resolved.json},
\texttt{plan.json}, \texttt{corpus-selected.jsonl}, \texttt{run.db},
\texttt{report.json} e os arquivos em \texttt{exports/}. Antes de uma
publicacao externa, recomenda-se revisar licencas do corpus, politicas de
redistribuicao e uma amostra manual dos exemplos aceitos.

\bibliographystyle{icml2026}
\begin{thebibliography}{9}

\bibitem[Kwon et~al.(2023)]{kwon2023vllm}
Woosuk Kwon, Zhuohan Li, Siyuan Zhuang, Ying Sheng, Lianmin Zheng, Cody Hao Yu,
Joseph E. Gonzalez, Hao Zhang, and Ion Stoica.
\newblock Efficient memory management for large language model serving with
PagedAttention.
\newblock In \emph{Proceedings of SOSP}, 2023.

\bibitem[Ouyang et~al.(2022)]{ouyang2022training}
Long Ouyang, Jeffrey Wu, Xu Jiang, Diogo Almeida, Carroll Wainwright, Pamela
Mishkin, Chong Zhang, Sandhini Agarwal, Katarina Slama, Alex Ray, et~al.
\newblock Training language models to follow instructions with human feedback.
\newblock In \emph{Advances in Neural Information Processing Systems}, 2022.

\bibitem[Patel et~al.(2026)]{patel2026fineinstructions}
Ajay Patel, Colin Raffel, and Chris Callison-Burch.
\newblock FineInstructions: Scaling Synthetic Instructions to Pre-Training
Scale.
\newblock \emph{arXiv preprint arXiv:2601.22146}, 2026.

\bibitem[Taori et~al.(2023)]{taori2023alpaca}
Rohan Taori, Ishaan Gulrajani, Tianyi Zhang, Yann Dubois, Xuechen Li, Carlos
Guestrin, Percy Liang, and Tatsunori B. Hashimoto.
\newblock Stanford Alpaca: An instruction-following LLaMA model.
\newblock Technical report, Stanford CRFM, 2023.

\bibitem[Wang et~al.(2023)]{wang2023selfinstruct}
Yizhong Wang, Yeganeh Kordi, Swaroop Mishra, Alisa Liu, Noah A. Smith, Daniel
Khashabi, and Hannaneh Hajishirzi.
\newblock Self-Instruct: Aligning language models with self-generated
instructions.
\newblock In \emph{Proceedings of ACL}, 2023.

\bibitem[Wei et~al.(2022)]{wei2022finetuned}
Jason Wei, Maarten Bosma, Vincent Y. Zhao, Kelvin Guu, Adams Wei Yu, Brian
Lester, Nan Du, Andrew M. Dai, and Quoc V. Le.
\newblock Finetuned language models are zero-shot learners.
\newblock In \emph{International Conference on Learning Representations}, 2022.

\end{thebibliography}

\end{document}
